\section{Inner product spaces}
An inner product assigns real numbers $\langle u,v\rangle$ and
satisfies symmetry, additivity, homogeneity, and positive definiteness.
The associated norm is $\|u\|=\sqrt{\langle u,u\rangle}$;
orthogonality means $\langle u,v\rangle=0$. The formula depends
on the space: a weighted coordinate product, evaluation of polynomials,
and integration of functions are different inner products.
For a nonzero vector $p$, projection onto its span is
$\langle q,p\rangle p/\langle p,p\rangle$.
For an orthogonal list, add these separate projections. Always use
the inner product stipulated in the exercise rather than substituting
the ordinary coordinate dot product.

\practice{6.7.P1}{A zero vector has zero inner product with every vector}
By homogeneity, $\langle0,v\rangle=\langle0u,v\rangle
=0\langle u,v\rangle=0$, for any $u$ in the space.
Symmetry then gives $\langle v,0\rangle=\langle0,v\rangle=0$.
The choice of $u$ does not matter, because its scalar multiple is zero.

\practice{6.7.P2}{Additivity in the second argument}
Use symmetry to move the sum into the first argument, where additivity
is an axiom:
\[
 \langle u,v+w\rangle=\langle v+w,u\rangle
 =\langle v,u\rangle+\langle w,u\rangle
 =\langle u,v\rangle+\langle u,w\rangle.
\]
The last equality applies symmetry to each term separately.

\exercise{6.7.1}{A weighted norm and its orthogonal line}
Example 1 uses $\langle u,v\rangle=4u_1v_1+5u_2v_2$.
\step{(a) Apply those weights to the stated vectors.}
For $x=(1,1)$ and $y=(5,-1)$,
$\|x\|^2=4+5=9$, $\|y\|^2=4(25)+5=105$,
and $\langle x,y\rangle=20-5=15$.
Thus $\|x\|=3$, $\|y\|=\sqrt{105}$, and
$|\langle x,y\rangle|^2=225$.
\step{(b) Solve the orthogonality equation.}
A vector $z=(z_1,z_2)$ is orthogonal to $y$ exactly when
$20z_1-5z_2=0$, or $z_2=4z_1$.
\answer{All such vectors are $t(1,4)$, $t\in\R$, including zero.}

\exercise{6.7.2}{Verify Cauchy--Schwarz with the weighted product}
For $x=(3,-2)$ and $y=(-2,1)$,
\[
 \langle x,y\rangle=4(3)(-2)+5(-2)(1)=-34,
\]
\[
 \|x\|^2=4(9)+5(4)=56,\qquad
 \|y\|^2=4(4)+5(1)=21.
\]
Therefore $|\langle x,y\rangle|^2=1156\le1176=\|x\|^2\|y\|^2$.
Both sides of $|\langle x,y\rangle|\le\|x\|\|y\|$ are nonnegative,
so the squared comparison proves the desired inequality.

\textbf{Evaluation convention for Exercises 3--8.}
The inner product is the dot product of the three value lists at
$-1,0,1$, not of the polynomial coefficient lists.
For Exercise 3, $p$ and $q$ have value vectors $(3,4,5)$ and
$(1,5,1)$. For Exercise 4, their value vectors are $(-4,0,2)$
and $(5,3,5)$.

\exercise{6.7.3}{Evaluate before multiplying}
For $p(t)=4+t$ and $q(t)=5-4t^2$,
\[
 \langle p,q\rangle=p(-1)q(-1)+p(0)q(0)+p(1)q(1)
 =3(1)+4(5)+5(1)=\boxed{28}.
\]
Each product uses values at the same input.

\exercise{6.7.4}{A negative inner product is allowed}
For $p(t)=3t-t^2$ and $q(t)=3+2t^2$,
\[
 \langle p,q\rangle=(-4)(5)+0(3)+2(5)=\boxed{-10}.
\]
Only self-products must be nonnegative; products of different vectors
can be negative, as this example shows.

\exercise{6.7.5}{Norms from the first pair of value lists}
Using the value vectors $(3,4,5)$ and $(1,5,1)$,
\[
 \|p\|=\sqrt{9+16+25}=\boxed{5\sqrt2},\qquad
 \|q\|=\sqrt{1+25+1}=\boxed{3\sqrt3}.
\]
The norm is the square root of the sum of squared values, not its square.

\exercise{6.7.6}{Norms from the second pair}
The value vectors are $(-4,0,2)$ and $(5,3,5)$, so
\[
 \|p\|=\sqrt{16+0+4}=\boxed{2\sqrt5},\qquad
 \|q\|=\sqrt{25+9+25}=\boxed{\sqrt{59}}.
\]
These are norms in the specified evaluation inner product.

\exercise{6.7.7}{Project onto the first polynomial line}
For the pair in Exercise 3, $\langle q,p\rangle=28$ and
$\langle p,p\rangle=50$. Since $p\ne0$, projection onto
$\Span\{p\}$ is
\[
 \operatorname{proj}_{\Span\{p\}}q
 =\frac{28}{50}p=\boxed{\frac{14}{25}(4+t)}.
\]
A check is $\langle q-(14/25)p,p\rangle=28-(14/25)50=0$.
Thus the residual is orthogonal to the whole one-dimensional target space.

\exercise{6.7.8}{Project onto the second polynomial line}
For the pair in Exercise 4, the coefficient is
$\langle q,p\rangle/\langle p,p\rangle=-10/20=-1/2$.
Hence the projection is
$\boxed{-\tfrac12(3t-t^2)=(t^2-3t)/2}$.
Its residual check is $-10-(-1/2)20=0$.
A negative projection coefficient is perfectly possible.

\exercise{6.7.9}{Remove the lower-degree part of a quadratic}
Use evaluation nodes $-3,-1,1,3$.
The value lists for $p_0,p_1,p_2$ are
$(1,1,1,1)$, $(-3,-1,1,3)$, and $(9,1,1,9)$.
\step{(a) Project onto the first two orthogonal polynomials.}
The first two lists have dot product zero. Also
$\langle p_2,p_0\rangle=20$, $\langle p_0,p_0\rangle=4$,
and $\langle p_2,p_1\rangle=-27-1+1+27=0$.
Thus the projection is $5p_0=5$.
\step{(b) Form and scale the orthogonal remainder.}
The remainder $t^2-5$ has values $(4,-4,-4,4)$.
Divide by four to get
$\boxed{q(t)=(t^2-5)/4}$, whose values are $(1,-1,-1,1)$.
It is orthogonal to both $p_0,p_1$ and is nonzero. The three
polynomials have distinct increasing degrees, so they are independent
and span the same space as $1,t,t^2$.

\exercise{6.7.10}{Best quadratic approximation to an odd cubic}
In the same four-node inner product, $t^3$ is orthogonal to the
even polynomials $p_0$ and $q$: the products cancel at opposite nodes.
Its only nonzero coefficient in the orthogonal list is on $p_1=t$:
\[
 \frac{\langle t^3,t\rangle}{\langle t,t\rangle}
 =\frac{81+1+1+81}{9+1+1+9}=\frac{41}{5}.
\]
\answer{The best approximation is $(41/5)t$. The best member of a
quadratic-polynomial space need not have a nonzero quadratic term.}

\exercise{6.7.11}{A cubic projected onto a five-node quadratic space}
Here evaluation is at $-2,-1,0,1,2$, and the source's orthogonal
polynomials are $p_0=1$, $p_1=t$, $p_2=t^2-2$.
Odd-even cancellation makes the projection coefficients on $p_0,p_2$
zero. The coefficient on $t$ is
\[
 \frac{\sum t_i^4}{\sum t_i^2}
 =\frac{16+1+0+1+16}{4+1+0+1+4}=\frac{17}{5}.
\]
\answer{The orthogonal projection is $(17/5)t$.}

\exercise{6.7.12}{Choose a scaled orthogonal cubic}
Subtract the projection from Exercise 11 to obtain
$r(t)=t^3-(17/5)t$. Its values at the five nodes are
$(-6/5,12/5,0,-12/5,6/5)$.
Multiplying by $5/6$ yields
\[
 \boxed{p_3(t)=(5t^3-17t)/6},
\]
with the requested values $(-1,2,0,-2,1)$. The remainder is
orthogonal to all of $\Span\{p_0,p_1,p_2\}$. Its degree is three,
so it is nonzero and adds a new direction. Thus the four polynomials
are an orthogonal basis for $\mathcal P_3$.

\exercise{6.7.13}{Transport the Euclidean product through an invertible map}
Define $\langle u,v\rangle=(Au)^T(Av)$.
Symmetry follows from symmetry of the dot product. For addition,
$A(u+w)=Au+Aw$, so
$\langle u+w,v\rangle=(Au+Aw)^TAv=\langle u,v\rangle+\langle w,v\rangle$.
Similarly $A(cu)=cAu$ gives homogeneity.
Finally $\langle u,u\rangle=\|Au\|^2\ge0$.
If it is zero, $Au=0$; invertibility gives $u=0$.
Conversely $u=0$ clearly gives zero. All four axioms hold.
The last step is why invertibility, rather than merely linearity of $A$,
was required.

\exercise{6.7.14}{The same construction for an injective transformation}
Define $\langle u,v\rangle=T(u)\cdot T(v)$.
Symmetry and the two linearity axioms follow exactly by applying
linearity of $T$ and the dot-product rules in $\R^n$.
The self-product is $\|T(u)\|^2$, nonnegative and zero exactly when
$T(u)=0$. Because $T$ is one-to-one and $T(0)=0$, this happens
only for $u=0$. Thus positive definiteness holds, proving an inner
product on $V$. Without injectivity, a nonzero kernel vector would
have self-product zero and violate that axiom.

\exercise{6.7.15}{Move a scalar out of the second argument}
Symmetry and homogeneity in the first argument give
\[
 \langle u,cv\rangle=\langle cv,u\rangle
 =c\langle v,u\rangle=c\langle u,v\rangle.
\]
These are real inner products, so no complex-conjugation convention is involved.

\exercise{6.7.16}{The distance between two orthonormal vectors}
Expand using bilinearity and symmetry:
\[
 \|u-v\|^2=\langle u,u\rangle-2\langle u,v\rangle+\langle v,v\rangle
 =1-0+1=2.
\]
Taking the nonnegative square root gives $\boxed{\|u-v\|=\sqrt2}$.

\exercise{6.7.17}{Recover an inner product from two norms}
Expand the two squared norms:
$\|u+v\|^2=\|u\|^2+2\langle u,v\rangle+\|v\|^2$ and
$\|u-v\|^2=\|u\|^2-2\langle u,v\rangle+\|v\|^2$.
Subtracting cancels both self-products and leaves
$\|u+v\|^2-\|u-v\|^2=4\langle u,v\rangle$.
Division by four proves the stated identity.

\exercise{6.7.18}{The parallelogram identity}
Add, instead of subtracting, the two expansions in Exercise 17.
The mixed terms $2\langle u,v\rangle$ and $-2\langle u,v\rangle$
cancel, leaving
\[
 \boxed{\|u+v\|^2+\|u-v\|^2=2\|u\|^2+2\|v\|^2}.
\]
The proof uses only inner-product axioms, so it is not restricted to coordinate vectors.

\exercise{6.7.19}{Arithmetic mean dominates geometric mean}
Use the ordinary dot product on the two supplied coordinate vectors.
Their dot product is $2\sqrt{ab}$, while each squared norm is
$a+b$. Cauchy--Schwarz gives
$2\sqrt{ab}\le\sqrt{a+b}\sqrt{a+b}=a+b$,
where nonnegativity of $a,b$ makes the square roots defined and the
left side nonnegative. Divide by two:
$\boxed{\sqrt{ab}\le(a+b)/2}$. The case $a=b=0$ is included.

\exercise{6.7.20}{The square of a mean versus the mean of squares}
For $u=(a,b)^T$ and $v=(1,1)^T$, squared Cauchy--Schwarz gives
$(a+b)^2\le(a^2+b^2)(1+1)=2(a^2+b^2)$.
Divide by four to obtain
\[
 \boxed{\left(\frac{a+b}{2}\right)^2\le\frac{a^2+b^2}{2}}.
\]
This holds for all real $a,b$, with no nonnegativity assumption needed.

\exercise{6.7.21}{An integral inner product that vanishes}
The product expands as
$(1-3t^2)(t-t^3)=t-4t^3+3t^5$.
Therefore
\[
 \langle f,g\rangle=\int_0^1(t-4t^3+3t^5)\,dt
 =\frac12-1+\frac12=\boxed0.
\]
Both functions are continuous, so this is the specified inner product;
they are orthogonal even though neither is zero.

\exercise{6.7.22}{Another orthogonal pair of functions}
Expanding $(5t-3)(t^3-t^2)$ gives $5t^4-8t^3+3t^2$.
Thus
$\langle f,g\rangle=\int_0^1(5t^4-8t^3+3t^2)\,dt
=1-2+1=\boxed0$.
The integral of a product can cancel even when the functions have nonzero norms.

\exercise{6.7.23}{An integral norm is a square root}
For $f(t)=1-3t^2$,
\[
 \|f\|^2=\int_0^1(1-6t^2+9t^4)\,dt
 =1-2+\frac95=\frac45.
\]
Hence $\boxed{\|f\|=2/\sqrt5}$. Nonnegativity determines the choice
of square root.

\exercise{6.7.24}{The norm of a cubic-minus-quadratic function}
For $g(t)=t^3-t^2$,
\[
 \|g\|^2=\int_0^1(t^6-2t^5+t^4)\,dt
 =\frac17-\frac13+\frac15=\frac1{105}.
\]
Thus $\boxed{\|g\|=1/\sqrt{105}}$.

\exercise{6.7.25}{An orthogonal polynomial basis on a symmetric interval}
On $[-1,1]$, oddness gives $\langle1,t\rangle=0$ and
$\langle t,t^2\rangle=0$. The projection of $t^2$ onto $1$
has coefficient $(\int_{-1}^1t^2)/(\int_{-1}^11)=(2/3)/2=1/3$.
Thus $t^2-1/3$ is orthogonal to both lower-degree polynomials.
A conventional scaling gives the basis
\[
 \boxed{\{1,\ t,\ (3t^2-1)/2\}}.
\]
The last polynomial has integral zero and its product with $t$ is odd.
All three are nonzero of distinct degrees, so they form a basis of
$\Span\{1,t,t^2\}$. Their squared norms are $2,2/3,2/5$, respectively.

\exercise{6.7.26}{Repeat the orthogonalization on a wider interval}
On $[-2,2]$, the first two polynomials are again orthogonal by oddness.
The coefficient to subtract from $t^2$ is
$(\int_{-2}^2t^2)/(\int_{-2}^21)=(16/3)/4=4/3$.
Scaling the remainder by $3$ gives
$\boxed{\{1,t,3t^2-4\}}$.
The last polynomial has integral $3(16/3)-4(4)=0$, and its
product with $t$ is odd. Distinct degrees give independence and the
same polynomial span. All squared norms are positive; the last is $256/5$.

\exercise{6.7.27}{Complete the five-node orthogonal polynomial basis}
Use the evaluation nodes $-2,-1,0,1,2$ and start with
$p_0=1,p_1=t,p_2=t^2-2$. They are orthogonal.
For $t^3$, only the projection onto $t$ is nonzero, with coefficient
$17/5$ as computed in Exercise 11. Set $r_3=t^3-(17/5)t$.
For $t^4$, oddness makes the projections onto $p_1$ and $r_3$
zero. The constant projection coefficient is $34/5$, while
\[
 \frac{\langle t^4,p_2\rangle}{\langle p_2,p_2\rangle}
 =\frac{(\sum t_i^6)-2(\sum t_i^4)}{14}
 =\frac{130-68}{14}=\frac{31}{7}.
\]
Hence the final remainder is
\[
 r_4=t^4-\frac{34}{5}-\frac{31}{7}(t^2-2)
 =t^4-\frac{31}{7}t^2+\frac{72}{35}.
\]
One convenient rescaled basis is
\[
 \boxed{\{1,t,t^2-2,5t^3-17t,35t^4-155t^2+72\}}.
\]
Their distinct degrees prove independence, and Gram--Schmidt preserves
the original span $\mathcal P_4$. Exact dot products of their five-node
value vectors are zero for every distinct pair, verifying orthogonality.

\exercise{6.7.28}{Orthogonalize powers of cosine by integration}
Start with $q_0=1$ and $q_1=\cos t$: their integral product over
$[0,2\pi]$ is zero. Since $\int\cos^2t=\pi$ and
$\int1=2\pi$, the projection of $\cos^2t$ onto the constant
is $1/2$. Its projection onto $\cos t$ is zero because
$\int_0^{2\pi}\cos^3t\,dt=0$. Thus
$q_2=\cos^2t-1/2$.
For $\cos^3t$, the constant and $q_2$ projections vanish by
cancellation over a half-period. The remaining coefficient is
\[
 \frac{\int_0^{2\pi}\cos^4t\,dt}{\int_0^{2\pi}\cos^2t\,dt}
 =\frac{3\pi/4}{\pi}=\frac34.
\]
Therefore $q_3=\cos^3t-(3/4)\cos t$, giving the orthogonal basis
\[
 \boxed{\{1,\cos t,\cos^2t-1/2,\cos^3t-(3/4)\cos t\}}.
\]
The squared norms are $2\pi,\pi,\pi/4,\pi/16$, all nonzero.
Equivalently the last two are $(1/2)\cos2t$ and $(1/4)\cos3t$,
so the basis spans the same four-dimensional space as the original list.
